% Satellite Constellation Argument-of-Latitude Optimisation
% 10-minute Beamer presentation sketch
% Build: make build/Project-Presentation.pdf

\documentclass[aspectratio=169, 10pt]{beamer}

\usetheme[lang=en,hr=false]{NewPwr}
\usepackage{amsmath, amssymb}
\usepackage{booktabs}
\usepackage{graphicx}
\usepackage{listings}
\lstset{basicstyle=\ttfamily\tiny, breaklines=false}
\usepackage{tikz}
\usetikzlibrary{arrows.meta, decorations.pathreplacing}

% ── Metadata ──────────────────────────────────────────────────────────────────
\title{Argument-of-Latitude Optimization\\in a Satellite Constellation}
\subtitle{Numerical Methods and Optimization --- Project}
\author{Sergiusz Warga 230757}
\institute{Wrocław University of Science and Technology\\
  Tutor: dr~inż.~Mateusz~Gabor\\
  Group: Wednesday 13:15}
\date{2026-06-10}

% ── Convenience macros ────────────────────────────────────────────────────────
\newcommand{\dv}{\Delta v}
\newcommand{\du}{\delta}

% ==============================================================================
\begin{document}

\maketitle

% ── Outline ───────────────────────────────────────────────────────────────────
\begin{frame}{Outline}
  \tableofcontents
\end{frame}

% ==============================================================================
\section{Motivation}

% ── 1. Motivation ─────────────────────────────────────────────────────────────
\begin{frame}{Motivation}
  \begin{columns}[T]
    \begin{column}{0.55\textwidth}
      \begin{itemize}[<+->]
        \item I don't have time. I asked our Mission Analysis
              team for a real problem I could implement for this course.
        \item Low-Earth-orbit constellations rely on
              \textbf{uniform angular spacing} to provide continuous coverage.
        \item Atmospheric drag, solar pressure, and Earth's oblateness slowly
              shift each satellite's angular position.
        \item Periodically, small \textbf{phasing manoeuvres} must restore equal
              spacing, but fuel is finite and unevenly distributed across the
              fleet.
        \item \textbf{Goal:} find the cheapest set of manoeuvres that re-spaces
              the constellation while keeping fuel levels as \emph{balanced} as
              possible.
      \end{itemize}
    \end{column}
    \begin{column}{0.42\textwidth}
      \only<1-5>{\begin{tikzpicture}[scale=1.15]
        \draw[thick, gray!60] (0,0) circle (1.4cm);
        \fill[Logo!80] (0,0) circle (0.12cm);
        % current positions only
        \foreach[count=\i] \a in {40,105,170,245,305,350}{
          \fill[Logo] ({1.4*cos(\a)},{1.4*sin(\a)}) circle (3.5pt)
            node[font=\tiny, above right] {$S_{\i}$};
        }
        \draw[->, gray!60] (0,0)--(1.6,0) node[right, font=\tiny]{$\Omega$};
      \end{tikzpicture}}
      \only<6->{\begin{tikzpicture}[scale=1.15]
        \draw[thick, gray!60] (0,0) circle (1.4cm);
        \fill[Logo!80] (0,0) circle (0.12cm);
        % ideal slots (dashed)
        \foreach \k in {0,...,5}{
          \draw[gray!30, dashed] (0,0) -- ({1.4*cos(60*\k)},{1.4*sin(60*\k)});
          \fill[gray!40] ({1.4*cos(60*\k)},{1.4*sin(60*\k)}) circle (2.5pt);
        }
        % current (perturbed) positions
        \foreach[count=\i] \a in {40,105,170,245,305,350}{
          \fill[Logo] ({1.4*cos(\a)},{1.4*sin(\a)}) circle (3.5pt)
            node[font=\tiny, above right] {$S_{\i}$};
        }
        \draw[->, gray!60] (0,0)--(1.6,0) node[right, font=\tiny]{$\Omega$};
        \node[font=\scriptsize, gray!60] at (0,-1.75) {grey: ideal \quad red: current};
      \end{tikzpicture}}
    \end{column}
  \end{columns}
\end{frame}

% ==============================================================================
\section{Orbital Mechanics Background}

% ── 3. Circular orbits and Kepler's 3rd law ───────────────────────────────────
\begin{frame}{Circular Orbits --- The Key Relationship}
  \begin{columns}[T]
    \begin{column}{0.55\textwidth}
      \onslide<1->{From Newton's law and centripetal acceleration, the orbital
      speed of a satellite at radius~$r$ is:
      \[
        v_c = \sqrt{\frac{\mu}{r}}, \qquad
        \mu = GM_\oplus \approx 3.986 \times 10^{14}\ \frac{\text{m}^3}{\text{s}^2}
      \]}
      \onslide<2->{The orbital period follows from Kepler's third law:
      \[
        T = \frac{2\pi r}{v_c} = 2\pi\sqrt{\frac{r^3}{\mu}}
      \]}
      \onslide<3->{\begin{block}{The key insight}
        \textbf{Lower orbit $\Rightarrow$ faster speed, shorter period.}\\
        \textbf{Higher orbit $\Rightarrow$ slower speed, longer period.}
      \end{block}}
    \end{column}
    \begin{column}{0.42\textwidth}
      \onslide<1->{\begin{tikzpicture}[scale=1.4]
        \draw[gray!50, thick] (0,0) circle (1cm);
        \fill[Logo!80] (0,0) circle (0.07cm)
              node[above right, font=\scriptsize]{$\oplus$};
        \draw[->, thick, Logo] (0,0) -- ({cos(55)},{sin(55)})
              node[midway, above left, font=\small]{$r$};
        \fill[Logo] ({cos(55)},{sin(55)}) circle (0.05cm)
              node[above right, font=\scriptsize]{$S$};
        \draw[->, thick, orange!80!black]
              ({cos(55)},{sin(55)}) -- ({cos(55)-0.35*sin(55)},{sin(55)+0.35*cos(55)})
              node[above, font=\scriptsize]{$v_c$};
        \draw[->, gray!70] (0.35,0) arc[start angle=0,end angle=55,radius=0.35cm]
              node[right, font=\scriptsize]{$u$};
        \draw[->, gray!60] (0,0)--(1.2,0) node[right, font=\tiny]{$\Omega$};
      \end{tikzpicture}\\[4pt]
      {\footnotesize The \textbf{argument of latitude}~$u$ is the angle from
      the ascending node $\Omega$ to the satellite's current position,
      measured along the orbit.}}
    \end{column}
  \end{columns}
\end{frame}

% ── 4. Orbit Phasing --- concept ─────────────────────────────────────────────
\begin{frame}{Orbit Phasing --- How It Works}
  \begin{columns}[T]
    \begin{column}{0.50\textwidth}
      \textbf{Need to advance (positive $\du$)?}
      \begin{enumerate}[<+->]
        \item Fire \textbf{retrograde} --- slow down.
        \item Satellite drops to a \emph{lower} orbit.
        \item Faster orbital speed $\Rightarrow$ shorter period.\\
              \textit{Completes one lap before the reference.}
        \item Fire \textbf{prograde} --- restore altitude.
      \end{enumerate}
      \pause\medskip
      \textbf{Need to fall back (negative $\du$)?}
      \begin{enumerate}[<+->]
        \item Fire \textbf{prograde} --- speed up.
        \item Satellite climbs to a \emph{higher} orbit.
        \item Slower speed $\Rightarrow$ longer period.\\
              \textit{Completes one lap behind the reference.}
        \item Fire \textbf{retrograde} --- restore altitude.
      \end{enumerate}
    \end{column}
    \begin{column}{0.46\textwidth}
      \onslide<10->{\begin{tikzpicture}[scale=1.5]
        % circular orbit (radius 1, Earth at origin)
        \draw[Logo, thick] (0,0) circle (1cm);
        % phasing orbit: apoapsis = 1 at (1,0), periapsis = 0.6 at (-0.6,0)
        % a=0.8, b=sqrt(0.64-0.04)=0.775, center=(0.2,0)
        \draw[red!60, thick, dashed] (0.2,0) ellipse (0.8cm and 0.775cm);
        % Earth
        \fill[Logo!80] (0,0) circle (0.06cm);
        % both burns happen at the same point: apoapsis (1,0)
        \fill[red!80] (1,0) circle (0.06cm);
        \node[font=\tiny, red!70, right] at (1.05,0.12) {burns 1 \& 2};
        % travel direction arrows
        \draw[->, Logo, thick] ({cos(100)},{sin(100)}) -- ({cos(101)},{sin(101)});
        \draw[->, red!50, thick]
              ({0.2+0.8*cos(100)},{0.775*sin(100)}) --
              ({0.2+0.8*cos(101)},{0.775*sin(101)});
        % labels
        \node[font=\tiny, Logo]    at (-0.3,  0.88) {circular orbit};
        \node[font=\tiny, red!60]  at (-0.35, -0.60) {phasing orbit};
      \end{tikzpicture}}
    \end{column}
  \end{columns}
\end{frame}

% ── 5. Delta-v formula ────────────────────────────────────────────────────────
\begin{frame}{Delta-$v$ for a Phasing Manoeuvre}
  \onslide<1->{To shift by angle~$\du$ in exactly one revolution, the phasing
  orbit semi-major axis is found from the period ratio:
  \[
    \frac{T_\text{ph}}{T_c} = 1 - \frac{\du}{2\pi}
    \;\Longrightarrow\;
    \frac{a_\text{ph}}{r} = \left(1 - \frac{\du}{2\pi}\right)^{2/3}
    \;\Longrightarrow\;
    \boxed{a_\text{ph} = \frac{r}{\bigl(1 + \du/(2\pi)\bigr)^{2/3}}}
  \]}
  \onslide<2->{The speed at the manoeuvre point on the phasing orbit
  (vis-viva equation):
  \[
    v_\text{ph} = \sqrt{\mu\!\left(\frac{2}{r} - \frac{1}{a_\text{ph}}\right)}
  \]}
  \onslide<3->{Both burns have equal magnitude, so the total cost is:
  \[
    \dv = 2\,\bigl|v_\text{ph} - v_c\bigr|
  \]}
  \begin{itemize}[<+->]
    \item $\du > 0$: $a_\text{ph} < r$, lower/faster orbit ---
          first burn is \textbf{retrograde}.
    \item $\du < 0$: $a_\text{ph} > r$, higher/slower orbit ---
          first burn is \textbf{prograde}.
  \end{itemize}
\end{frame}

% ── Problem Statement ─────────────────────────────────────────────────────────
\section{Problem Statement}

\begin{frame}{Problem Statement}
  \textbf{Inputs} (arbitrary, supplied by the caller):
  \begin{itemize}[<+->]
    \item $N$ satellites on a circular orbit at altitude $h$,
          gravitational parameter $\mu$.
    \item Current arguments of latitude $\mathbf{u} \in \mathbb{R}^N$.
    \item Remaining $\dv$ budgets $\mathbf{f} \in \mathbb{R}^N$.
  \end{itemize}

  \pause\medskip
  \textbf{Find:} angular transfers $\boldsymbol{\du}$ such that:
  \begin{enumerate}[<+->]
    \item Satellites are \textbf{equally spaced} ($360°/N$ apart) after the
          manoeuvre.
    \item \textbf{Total} $\dv$ consumed by the fleet is minimised.
    \item \textbf{Remaining} fuel budgets are as equal as possible.
  \end{enumerate}

  \pause\medskip
  \textbf{Simplifying assumptions:} circular, co-planar orbits; two-burn
  phasing manoeuvre per satellite; no perturbations during the transfer.

  \pause\medskip
  \textbf{Example used here:} $N=6$, $h = 550\ \text{km}$ (SSO),
  $\mathbf{u} = [40°,\ 105°,\ 170°,\ 245°,\ 305°,\ 350°]^\top$,
  $\mathbf{f} = [100,\ 95,\ 102,\ 98,\ 100,\ 97]^\top$ m/s.
\end{frame}

% ==============================================================================
\section{Mathematical Formulation}

% ── 6. 1-D reduction ──────────────────────────────────────────────────────────
\begin{frame}{Reducing to a 1-D Optimisation}
  \onslide<1->{After sorting satellites by current~$u$, the equal-spacing
  constraint fixes all target positions given a single \emph{reference phase}
  $u_\text{ref} \in [0,\ 2\pi)$:
  \[
    u_{i,\text{target}} = u_\text{ref} + (i-1)\cdot\frac{2\pi}{N},
    \quad i = 1,\dots,N.
  \]}
  \onslide<2->{The transfer angle for satellite~$i$:
  \[
    \du_i(u_\text{ref}) = \operatorname{wrap}_{[-\pi,\pi]}\!\bigl(
      u_{i,\text{target}} - u_i\bigr)
  \]}
  \onslide<3->{Wrapping to $[-\pi,\pi]$ selects the \textbf{shortest path}
  automatically.}

  \onslide<4->{\medskip
  The full $N$-satellite scheduling problem reduces to a
  \textbf{scalar search over $u_\text{ref}$}, solved with \texttt{fminbnd}.}
\end{frame}

% ── 7. Optimisation strategies ────────────────────────────────────────────────
\begin{frame}{Optimisation Criteria and Strategies}
  \onslide<1->{Define:
  \[
    J_1(u_\text{ref}) = \sum_{i=1}^N \dv_i, \qquad
    J_2(u_\text{ref}) = \max_i \bigl| r_i - \bar{r} \bigr|,
    \quad r_i = f_i - \dv_i \text{ (remaining fuel)}
  \]}

  \onslide<2->{\begin{block}{Weighted strategy (single \texttt{fminbnd} call)}
    \[
      \min_{u_\text{ref}}\ W_1 J_1(u_\text{ref}) + W_2 J_2(u_\text{ref})
    \]
    Trade-off between total cost and fuel balance controlled by $W_1, W_2$.
  \end{block}}

  \onslide<3->{\begin{block}{Lexicographic strategy (two \texttt{fminbnd} calls)}
    \textbf{Stage 1:} $\min J_1 \;\Rightarrow\; J_1^*$.\\
    \textbf{Stage 2:} $\min J_2$ subject to
    $J_1 \le J_1^*(1 + \varepsilon)$.\\
    Strictly prioritises fuel economy; balance is secondary.
  \end{block}}
\end{frame}

% ==============================================================================
\section{Implementation and Results}

% ── 9. Strategy Comparison ───────────────────────────────────────────────────
\begin{frame}{Strategy Comparison}
  \begin{columns}[T]
    \begin{column}{0.48\textwidth}
      \onslide<1->{\begin{block}{Weighted ($W_1 = 1,\ W_2 = 10$)}
        \begin{itemize}
          \item Single \texttt{fminbnd} call.
          \item $W_2 \gg W_1$: strongly penalises fuel imbalance.
          \item May consume slightly more total $\dv$ to equalise
                remaining budgets.
        \end{itemize}
      \end{block}}
    \end{column}
    \begin{column}{0.48\textwidth}
      \onslide<2->{\begin{block}{Lexicographic}
        \begin{itemize}
          \item Two sequential \texttt{fminbnd} calls.
          \item Stage~1: minimise total $\dv$ strictly.
          \item Stage~2: minimise imbalance within
                $\dv \le \dv^*(1 + \varepsilon)$.
        \end{itemize}
      \end{block}}
    \end{column}
  \end{columns}

  \onslide<3->{\medskip
  \begin{center}
    \includegraphics[height=0.45\textheight]{figures/comparison.pdf}
  \end{center}}
\end{frame}

% ── 10. Numerical Results: Weighted ──────────────────────────────────────────
\begin{frame}[fragile]{Numerical Results --- Weighted ($W_1=1,\ W_2=10$)}
  \lstinputlisting{out_weighted.m}
\end{frame}

% ── 11. Numerical Results: Lexicographic ─────────────────────────────────────
\begin{frame}[fragile]{Numerical Results --- Lexicographic}
  \lstinputlisting{out_lexicographic.m}
\end{frame}

% ── 12. Manoeuvre plans ───────────────────────────────────────────────────────
\begin{frame}{Manoeuvre Plans}
  \begin{columns}[c]
    \begin{column}{0.48\textwidth}
      \onslide<1->{\centering
      \includegraphics[width=\linewidth]{figures/manoeuvre_weighted.pdf}}
    \end{column}
    \begin{column}{0.48\textwidth}
      \onslide<2->{\centering
      \includegraphics[width=\linewidth]{figures/manoeuvre_lexicographic.pdf}}
    \end{column}
  \end{columns}
\end{frame}

% ── 13. Before / After ────────────────────────────────────────────────────────
\begin{frame}{Before and After}
  \begin{columns}[c]
    \begin{column}{0.48\textwidth}
      \onslide<1->{\centering
      \includegraphics[width=\linewidth]{figures/constellation_initial.pdf}}
    \end{column}
    \begin{column}{0.48\textwidth}
      \onslide<2->{\centering
      \includegraphics[width=\linewidth]{figures/constellation_target.pdf}}
    \end{column}
  \end{columns}
  \onslide<2->{\begin{center}
    \footnotesize After manoeuvres: spacing = $60^\circ$ between each adjacent pair.
  \end{center}}
\end{frame}

% ==============================================================================
\section{Discussion and Conclusions}

% ── 12. Limitations ──────────────────────────────────────────────────────────
\begin{frame}{Limitations and Possible Extensions}
  \begin{itemize}[<+->]
    \item \textbf{Multi-revolution phasing} --- each transfer is currently
          fixed to one phasing orbit. Allowing $k$ revolutions reduces $\dv$
          proportionally to $1/k$ at the cost of longer transfer time.
    \item \textbf{Fuel-budget feasibility} --- if the computed $\dv_i$
          exceeds a satellite's budget $f_i$, the solver returns negative
          remaining fuel without raising an error.
    \item \textbf{Global optimality} --- \texttt{fminbnd} finds a local
          minimum; multi-start sampling could verify global optimality for
          non-unimodal objectives.
    \item \textbf{Perturbation model} --- ignoring $J_2$, drag, and solar
          pressure leads to optimistic $\dv$ estimates.
    \item \textbf{Non-circular orbits} --- generalising to elliptic reference
          orbits would require a Hohmann transfer followed by a phasing arc.
  \end{itemize}
\end{frame}

% ── Thank you ─────────────────────────────────────────────────────────────────
\begin{frame}{}
  \begin{center}
    \Large\textbf{Thank you --- Questions?}\\[6pt]
    \normalsize\texttt{github.com/EdwardEisenhauer/AAE-NMaO}
  \end{center}
\end{frame}

\end{document}
